\begin{annexitem}{annex:p17.rat93-ess-equivalence}{thm:p17-rat93-ess-equivalence}{ESS, rare mutants and local superiority}
\begin{proof}
Write \(p=\sigma\), \(q=\mu\), and
\[
 a=p^{\mathsf T}Ap-q^{\mathsf T}Ap,\qquad
 b=p^{\mathsf T}Aq-q^{\mathsf T}Aq.
\]
The incumbent's advantage against \((1-\varepsilon)p+\varepsilon q\)
is \((1-\varepsilon)a+\varepsilon b\). If \(a>0\), this remains
positive for sufficiently small positive \(\varepsilon\). If
\(a=0,b>0\), it is positive for every positive \(\varepsilon\).
Conversely positivity for all sufficiently small positive
\(\varepsilon\) forces \(a\geq0\), and forces \(b>0\) when
\(a=0\). This proves the two-clause/rare-mutant equivalence. It
also proves the Nash implication and the strict-Nash sufficiency.

Here is the finite-dimensional argument supplying the uniform
neighborhood needed for local superiority. Suppose \(p\) is ESS.
Set \(v=p^{\mathsf T}Ap\),
\(d_i=v-(Ap)_i\geq0\), and
\(I=\{i:d_i=0\}\). Nash's support property gives
\(\operatorname{supp}(p)\subseteq I\).
Consider the closed cone
\[
 C_I=\{z:\sum_i z_i=0,\ z_i=0\ (i\notin I),
                \ z_i\geq0\ (i\notin\operatorname{supp}(p))\}.
\]
Every nonzero \(z\) in this cone has some sufficiently small
\(\alpha>0\) for which \(q=p+\alpha z\) is feasible and tied
against \(p\). The second ESS inequality gives
\[
 p^{\mathsf T}Aq-q^{\mathsf T}Aq=-\alpha^2z^{\mathsf T}Az>0.
\]
By compactness of \(C_I\cap\{z:\lVert z\rVert=1\}\), if this
set is nonempty there is \(c>0\) with
\(z^{\mathsf T}Az\leq-c\lVert z\rVert^2\) on the whole cone.
If the cone has no nonzero vector, all its vectors below are zero
and the same bound may be used with any \(c>0\).

For a nearby feasible \(y=p+h\), let
\(m=\sum_{i\notin I}h_i\geq0\). Put \(r_i=h_i\) outside
\(I\), \(r_i=-mp_i\) on \(\operatorname{supp}(p)\), and
\(r_i=0\) elsewhere. Then \(k=h-r\in C_I\) and
\(\lVert r\rVert\leq2m\); also \(m\leq\sqrt n\lVert h\rVert\).
Using the Euclidean operator norm, expansion gives
\(-h^{\mathsf T}Ah\geq c\lVert k\rVert^2
-2\lVert A\rVert\lVert k\rVert\lVert r\rVert
-\lVert A\rVert\lVert r\rVert^2\).
Since \(\lVert k\rVert\leq(1+2\sqrt n)\lVert h\rVert\),
we may take \(K=4\lVert A\rVert(1+3\sqrt n)\) to obtain
\[
 -h^{\mathsf T}Ah\geq c\lVert k\rVert^2-Km\lVert h\rVert.
\]
If \(I\) is not the whole index set, let
\(d=\min_{i\notin I}d_i>0\). Since
\(-h^{\mathsf T}Ap=\sum_{i\notin I}d_ih_i\geq dm\), we get
\[
 p^{\mathsf T}Ay-y^{\mathsf T}Ay
 =-h^{\mathsf T}Ap-h^{\mathsf T}Ah
 \geq (d-K\lVert h\rVert)m+c\lVert k\rVert^2.
\]
For sufficiently small \(\lVert h\rVert\) this is strictly
positive whenever \(h\ne0\): either \(m>0\), or \(m=0\) and
\(k=h\ne0\). If \(I\) is the whole set, \(m=r=0\) and the
quadratic bound alone proves the assertion. Thus ESS implies local
superiority, including boundary strategies.

Conversely suppose \(p\) is locally superior. For any \(q\ne p\)
let \(z=q-p\) and use nearby \(y=p+\alpha z\), \(\alpha>0\):
\[
 p^{\mathsf T}Ay-y^{\mathsf T}Ay
       =-\alpha z^{\mathsf T}Ap-\alpha^2z^{\mathsf T}Az>0.
\]
Letting \(\alpha\downarrow0\) gives \(z^{\mathsf T}Ap\leq0\).
If it is strictly negative, the first ESS clause holds. If it is
zero, the displayed inequality gives \(z^{\mathsf T}Az<0\), and
\(p^{\mathsf T}Aq-q^{\mathsf T}Aq=-z^{\mathsf T}Az>0\), giving
the second clause. The examples proved in the next two items show
failure of the converse implications involving Nash and strict Nash.
\end{proof}
\end{annexitem}

\begin{annexitem}{annex:p17.rat93-ess-restrictions}{thm:p17-rat93-ess-restrictions}{Weak dominance and support uniqueness}
\begin{proof}
If \(q\) weakly dominates \(p\), then
\(q^{\mathsf T}Ap\geq p^{\mathsf T}Ap\). ESS would first require
the reverse weak inequality, hence equality. Its second clause would
then require \(p^{\mathsf T}Aq>q^{\mathsf T}Aq\), contrary to
dominance evaluated at \(q\). The dominator differs from the dominated
strategy because the definition contains a strict comparison somewhere.

Suppose \(p,q\) are distinct ESSs with identical support. Both are
Nash. All pure strategies on \(p\)'s support tie against \(p\), so
\(q^{\mathsf T}Ap=p^{\mathsf T}Ap\). The second ESS clause for
\(p\) requires \(p^{\mathsf T}Aq>q^{\mathsf T}Aq\). But every
pure strategy on \(q\)'s support ties against \(q\); since
\(p\) has that support, these two quantities are equal. This is a
contradiction. There are \(2^n-1\) nonempty supports. The RPS item
below gives a game with no ESS.
\end{proof}
\end{annexitem}

\begin{annexitem}{annex:p17.rat93-static-calculations}{thm:p17-rat93-static-calculations}{Invasion, Hawk--Dove and declining mean welfare}
\begin{proof}
For \(A=\left(\begin{smallmatrix}2&1\\1&1\end{smallmatrix}\right)\),
let \(q\) be the opponent's \(S_1\) probability. Pure payoffs are
\(1+q\) and \(1\). Thus \(S_1\) is strictly best whenever
\(q>0\); at \(q=0\) all mixtures tie. The only symmetric Nash
strategies are \(S_1,S_2\). The former is strict against itself.
For the latter, a pure \(S_1\) mutant ties against \(S_2\) but
\(\pi(S_2\mid S_1)=1<2=\pi(S_1\mid S_1)\), so it is not ESS.
At mutant fraction \(\varepsilon\), those same pure payoffs differ
by \(\varepsilon\).

For Hawk--Dove, the pure payoff difference against Hawk share \(q\) is
\[
 \pi(H\mid q)-\pi(D\mid q)=\frac{V-Cq}{2}.
\]
The best response is Hawk for \(q<V/C\), Dove for \(q>V/C\), and
any mixture at equality. Because \(0<V/C<1\), pure profiles
\((H,D),(D,H)\) are Nash and the same-action profiles are not.
If either player mixes nontrivially, its opponent's Hawk share must
equal \(V/C\), which forces both to mix at that share. This gives
exactly the stated three profiles.
At \(q^*=V/C\), both pure payoffs tie, so every mutant mixture
also ties against \(\sigma^*\). Against mutant share \(q\),
the incumbent-minus-mutant payoff is
\[
 (q^*-q)\{\pi(H\mid q)-\pi(D\mid q)\}
       =(q^*-q)\frac{C(q^*-q)}2
       =\frac C2(q-q^*)^2.
\]
It is strictly positive for every different mutant. Thus the mixture
is ESS but not strict Nash.

For the welfare game, pure payoffs against \(S_1\) share \(q\)
are \(100-99q\) and \(100-100q\), with difference \(q\).
This proves weak domination and the same two symmetric Nash states.
\(S_1\) is strict against itself, hence ESS; \(S_2\) fails the
tie comparison since \(0<1\) against \(S_1\).
Every entry of each player's payoff matrix is at most \(100\), so
the profile giving \((100,100)\) cannot be Pareto dominated,
including by randomization. This is a game with supplied comparable
payoff numbers for its stated comparison, not a claim that arbitrary
utilities from different people may always be added.

At population share \(x\), replicator gives
\(\dot x=x(1-x)(W_1-W_2)=x^2(1-x)\). Its interior solutions
increase and can converge only to \(1\): a limit strictly between
zero and one would have a positive derivative bounded away from zero.
Mean payoff is
\[
 \phi=x(100-99x)+(1-x)(100-100x)=100-100x+x^2.
\]
Along an interior trajectory,
\(\dot\phi=(-100+2x)x^2(1-x)<0\).
The stable invasion-resistant endpoint has self-play payoff \(1\),
whereas the other Nash endpoint gives \(100\) to both players.
\end{proof}
\end{annexitem}

\begin{annexitem}{annex:p17.rat93-protocols}{thm:p17-rat93-protocols}{Mean dynamics, support and stationary incentives}
\begin{proof}
Substituting imitation into the mean equation gives
\[
 \dot p_i=p_i\sum_jp_j\{[W_i-W_j]_+-[W_j-W_i]_+\}
          =p_i\sum_jp_j(W_i-W_j)=p_i(W_i-\phi),
\]
using \([a]_+-[-a]_+=a\). BNN's inflow is
\(\sum_jp_j[W_i-\phi]_+=[W_i-\phi]_+\); its outflow is
the stated multiple of \(p_i\). Smith substitution gives its two
displayed sums directly.

Summing any mean equation cancels each directed inflow with its
outflow. At \(p_i=0\) the remaining inflow is nonnegative. More
explicitly, writing that coordinate equation as
\(\dot p_i=a_i(t)-b_i(t)p_i\), with nonnegative inflow
\(a_i\) and total switch rate \(b_i\), the integrating-factor
formula preserves nonnegativity. Together with the conserved sum
this proves simplex invariance. Local Lipschitzness gives local
existence and uniqueness; compactness bounds the vector field and
extends each solution for every finite forward time. For replicator,
\[
 p_i(t)=p_i(0)\exp\!\left(\int_0^t[W_i(p(s))-\phi(p(s))]\,ds\right).
\]
The bounded integrand makes the finite-time support assertions immediate.
For BNN a missing strategy with \(W_i>\phi\) has positive inflow.
For Smith a missing strategy strictly better than some represented
strategy has positive inflow. The bad-RPS vertex in the body witnesses
both possibilities.

For BNN set \(\beta_i=[W_i-\phi]_+\) and \(B=\sum_i\beta_i\).
At rest \(\beta_i=p_iB\). If \(B>0\), every represented strategy
has \(\beta_i>0\), hence \(W_i>\phi\). Their weighted average
cannot then be \(\phi\). Therefore \(B=0\), all
\(W_i\leq\phi\), and the state is Nash. Conversely Nash implies
all \(\beta_i=0\), giving rest.

For Smith, multiplying the mean equation by payoffs and summing gives
\[
 W^{\mathsf T}\dot p
  =\sum_{i,j}p_i[W_j-W_i]_+(W_j-W_i)
  =\sum_{i,j}p_i[W_j-W_i]_+^2.
\]
At rest this is zero, so no represented strategy has an improving
alternative. Thus all represented strategies attain the global maximum
payoff, the Nash condition. Conversely at Nash, every positive-mass
incumbent has no higher-payoff destination, so every nonzero transition
flow vanishes. Replicator's simplex vertices are always rest points;
the bad-RPS vertex is not Nash since Paper is strictly better there.
\end{proof}
\end{annexitem}

\begin{annexitem}{annex:p17.rat93-replicator-dominance}{thm:p17-rat93-replicator-dominance}{Mixed domination and surviving weakly dominated shares}
\begin{proof}
Let \(\mu\) strictly dominate \(S_i\). The gap between their
expected payoffs is linear in the opponent distribution and strictly
positive on the compact simplex. Its minimum is some \(\eta>0\).
All coordinates remain positive at finite times, so
\[
 Z(t)=\frac{\prod_jp_j(t)^{\mu_j}}{p_i(t)}
\]
is defined, taking zero exponent factors as \(1\). Its logarithm has
derivative
\[
 \frac{d}{dt}\log Z
     =\sum_j\mu_j(W_j-\phi)-(W_i-\phi)
     =\pi(\mu\mid p)-\pi(S_i\mid p)\geq\eta.
\]
Hence \(Z(t)\geq Z(0)e^{\eta t}\). The numerator is at most
one, so \(p_i(t)\leq Z(0)^{-1}e^{-\eta t}\to0\).
This proof requires positivity of the coordinates used by \(\mu\),
as guaranteed by the source's full-interior hypothesis. If a strictly
dominating Defect strategy is absent in an all-Cooperate population,
the support invariance instead leaves that dominated vertex fixed.

For Figure~9.3.4 put \((x,y,z)=(p_1,p_2,p_3)\).
The pure payoffs are \(1,1-z,0\), proving the stated weak
dominance. At an interior state, ratio derivatives give
\[
 \frac{d}{dt}\log(z/x)=-1,\qquad
 \frac{d}{dt}\log(y/x)=-z.
\]
Thus \(z/x=(z(0)/x(0))e^{-t}\), and since \(x\leq1\),
\(0<z\leq(z(0)/x(0))e^{-t}\). In particular \(z\to0\) and
\(\int_0^\infty z(t)\,dt<\infty\). The second ratio therefore
converges to a finite strictly positive constant
\[
 L=\frac{y(0)}{x(0)}\exp\!\left(-\int_0^\infty z(t)\,dt\right).
\]
Since \(x+y=1-z\), it follows that
\((x,y,z)\to(1/(1+L),L/(1+L),0)\). Weakly dominated
\(S_2\) survives with positive share.
\end{proof}
\end{annexitem}

\begin{annexitem}{annex:p17.rat93-stability-bridges}{thm:p17-rat93-stability-bridges}{Nash, Lyapunov stability and the ESS Lyapunov function}
\begin{proof}
At a population Nash state all represented pure strategies attain the
mean payoff. Their replicator derivatives vanish, and the derivatives
of absent coordinates vanish by the population factor. Thus Nash
implies rest.

Let \(p^*\) be a Lyapunov-stable rest point that is not Nash.
Some \(j\) has \(W_j(p^*)-\phi(p^*)>0\); continuity gives
a neighborhood in which this difference is at least \(\eta>0\).
Stability would allow an initial state in this neighborhood with
\(p_j(0)>0\), arbitrarily close to \(p^*\), whose entire trajectory
remains in it. But then \(p_j(t)\geq p_j(0)e^{\eta t}\), eventually
exceeding one. This contradiction proves the stable-rest-to-Nash direction.

For the singleton simplex the stability and attraction assertions are
immediate. Otherwise suppose now that \(p^*\) is ESS, hence locally
superior by the first annex item. In a sufficiently small neighborhood every coordinate
on \(S^*=\operatorname{supp}(p^*)\) is positive. Define
\[
 V(p)=\sum_{i\in S^*}p_i^*\log\frac{p_i^*}{p_i}.
\]
Concavity of the logarithm gives
\[
 -V(p)=\sum_{i\in S^*}p_i^*\log(p_i/p_i^*)
       \leq\log\sum_{i\in S^*}p_i\leq0.
\]
Equality requires both total mass one on that support and constant
ratios \(p_i/p_i^*\), hence \(p=p^*\). Thus \(V\) is continuous,
nonnegative, and has a strict local minimum at \(p^*\).
Along replicator trajectories,
\[
 \dot V=-\sum_{i\in S^*}p_i^*(W_i-\phi)
         =p^{\mathsf T}Ap-(p^*)^{\mathsf T}Ap<0
                       \quad(p\ne p^*\text{ nearby}).
\]
For completeness, choose a small closed ball relative to the simplex
contained in this neighborhood. The positive minimum of \(V\) on
its boundary gives a sublevel neighborhood that trajectories cannot
leave, proving Lyapunov stability. Inside any smaller compact annulus
about \(p^*\), \(\dot V\) has a strictly negative maximum.
Consequently a trajectory in the invariant sublevel set spends only
finite total time in that annulus. Its bounded speed means that
infinitely many excursions a fixed distance from \(p^*\) would each
take a fixed positive time through such an annulus. That is impossible.
Thus the trajectory converges to \(p^*\), proving local attraction
and asymptotic stability. Boundary ESSs are included because only
their positive-support coordinates occur in \(V\).
\end{proof}
\end{annexitem}

\begin{annexitem}{annex:p17.rat93-rps-calculation}{thm:p17-rat93-rps-calculation}{Zero-sum RPS, no ESS and neutral cycles}
\begin{proof}
The RPS matrix is skew symmetric, so \(p^{\mathsf T}Ap=0\).
Its uniform strategy gives zero against every opponent, and every
strategy gives zero against the uniform strategy. Therefore any
Nash profile must give value zero. Its row-player deviation inequalities
give \(Aq\leq0\) componentwise; the components sum to zero, hence
\(Aq=0\). Solving gives \(q_1=q_2=q_3=1/3\). The column inequalities
give the same conclusion for the row mixture. Thus the Nash profile
is unique. Against the uniform incumbent every mutant ties, and
\(\pi(p^*\mid q)=\pi(q\mid q)=0\), violating the strict ESS
tie comparison for every \(q\ne p^*\).

The replicator equations are
\[
 \dot p_1=p_1(p_3-p_2),\quad
 \dot p_2=p_2(p_1-p_3),\quad
 \dot p_3=p_3(p_2-p_1).
\]
In the interior,
\(d\log(p_1p_2p_3)/dt=0\). By the arithmetic--geometric mean
inequality the product has its unique maximum \(1/27\) at \(p^*\).
For any prescribed neighborhood, the product's maximum on the compact
set outside it is strictly smaller. All sufficiently close initial
states have product above that smaller maximum and cannot leave the
neighborhood. This proves Lyapunov stability.

Any different interior state has constant product less than \(1/27\),
so cannot converge to \(p^*\). The interior level sets are compact
and avoid the boundary and the center. They are regular one-dimensional
closed curves: the constrained gradient of \(\log(p_1p_2p_3)\)
vanishes only at the center. The vector field is tangent to them and
never zero there, so each connected level curve is traversed periodically.
All vertices are rest points by support invariance, but a missing
strategy beats the resident at each vertex, making them non-Nash.
Thus dynamic neutral stability is neither ESS nor attraction.
\end{proof}
\end{annexitem}

\begin{annexitem}{annex:p17.rat93-regular-ess}{thm:p17-rat93-regular-ess}{Tangent definiteness and a nonregular ESS}
\begin{proof}
At a regular Nash \(p\), all tied best responses have support
contained in \(\operatorname{supp}(p)\), because every unused
strategy is strictly worse. For any different tied \(q\),
\(z=q-p\) is a nonzero zero-sum vector supported on that face,
and \(z^{\mathsf T}Ap=0\). Hence
\[
 p^{\mathsf T}Aq-q^{\mathsf T}Aq=-z^{\mathsf T}Az>0.
\]
Untied alternatives already satisfy the first ESS clause. This proves
regular-ESS sufficiency.

For \(A=\left(\begin{smallmatrix}1&1\\1&0\end{smallmatrix}\right)\),
take \(p=S_1\) and \(q=(1-a,a)\), \(0<a\leq1\).
Every strategy ties against \(p\), while
\(\pi(p\mid q)=1\) and \(\pi(q\mid q)=1-a^2\).
The second ESS difference is \(a^2>0\). But the unused strategy
\(S_2\) ties against \(S_1\), so quasistrictness fails.
For a two-strategy interior support, a nonzero tangent vector has
both coordinates nonzero; requiring \emph{proper} support containment
would exclude all such tests, explaining why that reading of the
source's support symbol is inadmissible.
\end{proof}
\end{annexitem}

\begin{annexitem}{annex:p17.rat93-two-populations}{thm:p17-rat93-two-populations}{The coupled Battle-of-the-Sexes equations}
\begin{proof}
For a two-strategy population, subtracting its two payoffs reduces
replicator to \(\dot x=x(1-x)(W_1-W_2)\). Row's difference is
\(y-2(1-y)=3y-2\); Column's is \(2x-(1-x)=3x-1\).
This proves the two equations and the source's role-specific means.

If \(0<x,y<1\), both derivatives vanish exactly at
\((x,y)=(1/3,2/3)\). On each edge, the remaining coordinate
can be at rest only at its endpoints, since its payoff difference
is nonzero. Thus the corners and the one interior point exhaust rest.
The best-response threshold is \(2/3\) for Row and \(1/3\) for
Column. The matched corners and the mixed point are mutual best
responses; at each mismatched corner some player has a strict improvement.
There are no additional Nash profiles because any nontrivial mixing
forces those thresholds for both opponents.

Near \((0,0)\), in the invariant rectangle
\(x<1/6,y<1/3\), both coordinates decrease, with rates bounded
above by negative constants times their own shares. They tend
exponentially to zero; small initial coordinates remain small, proving
asymptotic stability. Near \((1,1)\), use
\(x>2/3,y>5/6\): both deficits decrease exponentially and remain
small. The cross corners are unstable: on the edge \(y=0\),
any \(x<1\) moves away from \((1,0)\), while on \(x=0\), any
\(y<1\) moves away from \((0,1)\).

At \((1/3,2/3)\) the Jacobian is
\(\left(\begin{smallmatrix}0&2/3\\2/3&0\end{smallmatrix}\right)\),
whose eigenvectors \((1,1),(1,-1)\) have eigenvalues
\(2/3,-2/3\). The equilibrium is hyperbolic, so the usual local
linearization/stable-manifold characterization gives a saddle.
In particular, instability can be seen without linearization:
points with both coordinates just above their thresholds have both
coordinates increasing and move out of a fixed small rectangle.
This is a two-role stationary/stability calculation, not a one-population
ESS identification.
\end{proof}
\end{annexitem}

\begin{annexitem}{annex:p17.rat93-small-basin}{thm:p17-rat93-small-basin}{A strict ESS with an arbitrarily thin basin}
\begin{proof}
At the all-\(S_3\) state, payoffs to \(S_1,S_2,S_3\) are
\(\varepsilon,0,2\varepsilon\). Thus \(S_3\) is strict Nash,
and is a regular ESS because its face has no nonzero tangent vector.
The proved replicator bridge gives asymptotic stability.

Write \((x,y,z)=(p_1,p_2,p_3)\). The pure payoffs are
\(y+\varepsilon z,y,2\varepsilon z\), and direct averaging gives
\[
 \dot y=yz(y-\varepsilon x-2\varepsilon z).
\]
When \(y>2\varepsilon\), the bracket is positive because
\(\varepsilon x+2\varepsilon z\leq2\varepsilon\).
Thus \(y\) is nondecreasing there (constant when \(z=0\))
and cannot cross back below its initial positive value. Such a
trajectory cannot converge to \(S_3\), where \(y=0\).
Therefore the basin is inside \(y\leq2\varepsilon\).

In coordinates \((x,y)\), simplex area is \(1/2\). The strip
\(0\leq y\leq a\) has area
\(\int_0^a(1-y)\,dy=a-a^2/2\). Its relative area is
\(2a-a^2\), giving \(4\varepsilon-4\varepsilon^2\) for
\(a=2\varepsilon\). The intrinsic simplex area has the same
ratio because its coordinate Jacobian is constant. This is an upper
bound, not an assertion that every point in the strip belongs to
the basin. Nor is it a probability bound without a specified
initial-condition distribution.
\end{proof}
\end{annexitem}

\begin{annexitem}{annex:p17.rat93-global-unstable}{thm:p17-rat93-global-unstable}{Interior convergence with large excursions}
\begin{proof}
Write \((x,y,z)=(p_1,p_2,p_3)\). The pure payoffs are
\(2z,x,z\), and \(\phi=2xz+xy+z^2\).
If \(z>0\), \(S_3\) pays less than \(S_1\) and cannot have positive
share in a Nash state. Thus a Nash state must have \(z=0\). If then
\(x>0\), \(S_1\) pays less than \(S_2\) and cannot be represented
either. The only Nash state is therefore \(S_2\). All strategies tie
against it, while \(\pi(S_2\mid S_3)=0<1=\pi(S_3\mid S_3)\),
so it is not ESS.
For an initially interior solution all three shares remain positive.
The ratio derivatives are
\[
 \frac{d}{dt}\log(x/z)=z,\qquad
 \frac{d}{dt}\log(y/x)=x-2z.
\]
Put \(r=x/z\). If \(\int_0^\infty z\,dt<\infty\), the first
identity gives a finite positive limit for \(r\). An integrable
nonnegative function with bounded derivative tends to zero: otherwise
infinitely many peaks above some fixed height would, by the derivative
bound, each contribute a fixed positive amount to its integral.
Polynomial replicator flow has bounded derivatives on the simplex,
so \(z\to0\). Then \(x=rz\to0\) and \(y\to1\).

If instead \(\int_0^\infty z\,dt=\infty\), then \(r\to\infty\).
For all sufficiently late times \(r\geq4\), and hence
\(x-2z\geq x/2\). If the late-time integral of \(x\) is finite,
the same bounded-derivative argument gives \(x\to0\), and
\(z=x/r\to0\). If its integral is infinite, the second ratio
identity gives \(y/x\to\infty\). Since \(y\leq1\), this again
forces \(x\to0\), hence \(z\to0\) and \(y\to1\).
This exhausts the cases and proves convergence for every interior
initial state.

To disprove Lyapunov stability, start on the invariant boundary
\(x=0\), with \(z(0)>0\) arbitrarily small and
\(y(0)=1-z(0)\). There
\(\dot z=z^2(1-z)>0\), so the trajectory eventually reaches
\(z=1/2\), a fixed positive distance from \(S_2\).
These initial states approach \(S_2\) as \(z(0)\downarrow0\).
For each such trajectory, the time to reach \(z=1/2\) is finite.
Continuous dependence on initial conditions then allows a sufficiently
small positive \(x(0)\), with \(y(0)\) adjusted, to retain an
excursion to \(z>1/3\) before that finite time. These are interior
initial states arbitrarily close to \(S_2\), and by the first part
they eventually return and converge to it. Thus even full-interior
attraction does not supply Lyapunov stability or asymptotic stability.
\end{proof}
\end{annexitem}
